\documentclass[11pt,a4paper]{article}

\usepackage[margin=2.5cm]{geometry}
\usepackage{graphicx}
\usepackage{amsmath}

\begin{document}

\section*{Rough notes on the project}

In the following, I will give a brief overview of the progress with the first positive results and space for improvement.
The project is about learning a placement policy for serverless/edge workloads using a graph neural network (GNN) that is trained inside of a HeROsim simulator.
The simulator gives me many examples of ``this batch of requests + this system state $\rightarrow$ total RTT for a given placement'' which is the input for our GNN training pipeline.

\section*{Co-simulation}

For each sampled situation (infrastructure + current queues + a small batch of incoming requests) I enumerate all valid task--to--replica placements, and let the simulator run them to completion.
For every placement I record the resulting total RTT; across all placements for that situation I also keep track of the best RTT I have seen so far.
This gives me a dataset where each example is ``graph of the state + one concrete placement $\rightarrow$ measured RTT'', and from that I extract:
\begin{itemize}
  \item a \emph{label} for each task (which replica was part of the best placement),
  \item a big lookup table that maps full placements to RTT, which I later use in the regret loss.
\end{itemize}
Because the placements are allowed to be non-unique, multiple tasks can end up on the same replica, just like in the real system.

\section*{What the GNN sees}

For each small batch of requests I build a bipartite graph.
One side contains the tasks, the other side contains all possible compute replicas.
Edges connect a task to the replicas that are reachable over the network and can actually run that task.

Each node and edge has a small feature vector:
\begin{itemize}
  \item tasks know their type and which client they came from,
  \item replicas know their hardware type, whether they already host warm replicas, their queue length and some simple consolidation statistics,
  \item edges store an estimate of execution cost, network latency and some communication cost.
\end{itemize}
This way, the graph encodes both ``who acan execute whom'' and ``how expensive would it roughly be''.
The GNN processes this graph and produces a score for every (task, replica) edge; for each task I pick the replica with the highest score.

\section*{GNN architecture in a bit more detail}

Very roughly, the model has three parts:
\begin{itemize}
  \item a small MLP that turns the task features into task embeddings,
  \item another MLP that turns the replica features into replica embeddings,
  \item a graph neural network layer stack plus an edge MLP that use the graph structure to refine these embeddings and score each possible (task, replica) pair.
\end{itemize}

Concretely, tasks and replicas start as different feature vectors, but both are projected into the same embedding dimension.
These embeddings are stacked into a single node matrix and passed through several message-passing layers (a GIN-style architecture).
Message passing lets a task ``see'' not only its own features, but also information coming from nearby replicas and other tasks via the shared infrastructure.

After message passing, I split the node matrix back into task nodes and replica nodes.
For every edge in the graph I concatenate:
\begin{itemize}
  \item the final embedding of the task at that end of the edge,
  \item the final embedding of the replica at the other end,
  \item the edge feature vector (execution estimate, latency, warm/cold flag, communication cost).
\end{itemize}
This concatenated vector goes through another MLP (the edge scorer) which outputs a single scalar score for that edge.
Grouping these scores per task gives me a variable-length logit vector: one logit per feasible replica.
Because I work with non-unique placements, there is no constraint that different tasks must pick different replicas; each task just chooses the replica with the highest score among its candidates.

\section*{Training idea}

The training mixes two objectives:
\begin{itemize}
  \item a normal cross-entropy loss that says ``pick the same replica as the best placement the simulator found'',
  \item a regret-based loss that compares a good placement with a worse one and pushes the score gap to be at least as big as the RTT difference.
\end{itemize}
This way the model does not just copy labels, it also learns to strongly dislike obviously bad placements.

\section*{Workloads}

In the plots I use two simple workload labels that come from how I configure the trace generator:
\begin{itemize}
  \item \textbf{100-100}: about 100 requests per second for 100 seconds.
  \item \textbf{100-50}: about 100 requests per second for 50 seconds.
\end{itemize}
100-50 has 100k requests, while 100-100 has twice that many. That's because the amount of requests is calculate by multiplying rps * duration * number of clients (20).

\section*{RTT comparison}

In the simulator I compare several policies:
the GNN policy greedy (named knative), random and always-offload).
The plot shows the total request-response time (RTT) for different policies and for the two workloads (100-100 and 100-50).
Lower bars are better - smaller RTT.

\begin{figure}[h]
  \centering
  \includegraphics[width=\textwidth]{images/policy_comparison.png}
  \caption{Total RTT of different scheduling policies for the 100-100 and 100-50 workloads.}
\end{figure}

In both workloads the GNN policy gives the lowest total RTT.
The knative baseline comes next, and the random and always-offload style policies are clearly worse.

\end{document}